\documentclass[11pt,a4paper]{article}

% --- PACKAGES ---
\usepackage[utf8]{inputenc}
\usepackage[T1]{fontenc}
\usepackage[english]{babel}
\usepackage{geometry}
\geometry{a4paper, margin=1in}
\usepackage{graphicx}
\usepackage{booktabs}
\usepackage{amsmath}
\usepackage{amssymb}
\usepackage{hyperref}
\usepackage{charter} % Premium font look
\usepackage{microtype}
\usepackage{cite}
\usepackage{float}
\usepackage{caption}

% --- METADATA ---
\title{Developing Robotic Vision Robustness: Bridging the Sim-to-Real Gap through Bayesian Domain Randomization}
\author{Your Name \\ [p12289lp] \\ University of XXX}
\date{\today}

\begin{document}

\maketitle

\begin{abstract}
Deep Neural Networks (DNNs) frequently fail to generalize from simulation to real-world environments due to differences in visual fidelity and sensor noise. This report investigates the use of Bayesian Optimization and Domain Randomization to improve the robustness of Convolutional Neural Networks (CNNs) against Gaussian noise and blur. We evaluate a custom CNN model on the CIFAR-100 dataset against various levels of synthetic robotic sensor degradation. Our results demonstrate that while robust training incurs a minor ``Alignment Tax'' on pristine data, it prevents the catastrophic performance collapse observed in standard optimized baselines.
\end{abstract}

\section{Part A: Literature Review}
\label{sec:lit_review}

\subsection{Introduction: The Sim-to-Real Challenge}
Deep learning has revolutionized robotic perception, yet most models trained in static, clean environments are structurally brittle. In safety-critical applications such as autonomous navigation or robotic surgery, the inability to handle sensor noise or environmental shifts can lead to catastrophic failure.

\subsection{State of the Art in Robustness}
Early research into robotic transfer identified that simulations cannot perfectly model reality. Tobin et al. \cite{tobin2017domain} proposed \textit{Domain Randomization} as a solution, suggesting that by varying simulation parameters globally, a model can learn the underlying structural invariant features of a scene.

Parallel to this, Hendrycks and Dietterich \cite{hendrycks2019benchmarking} provided a standardized suite of image corruptions, highlighting that standard ImageNet-grade CNNs are hypersensitive to Gaussian noise. Geirhos et al. \cite{geirhos2018imagenet} expanded on this by proving that CNNs are heavily biased toward texture over global shape, which explains their vulnerability to surface-level noise like ISO grain or lens dirt.

\subsection{Regularization and Generalization}
Methods such as Random Erasing \cite{zhong2020random} and large-scale data collection \cite{levine2018learning} have been used to improve generalization. However, as the literature suggests, robustness is not a byproduct of dataset size alone but is often a result of explicit, adversarial-style regularization during the training phase.

\section{Part B: Experimental Setup}
\label{sec:methodology}

\subsection{Methodology and Dynamic Noise Injection}
We utilize CIFAR-100 as a proxy for complex object classification. To simulate robotic sensor degradation, we implemented a custom \texttt{tf.data} pipeline that injects Gaussian noise at various severities ($0.1$ to $0.4$).

\subsection{Bayesian Hyperparameter search}
We leveraged the Optuna framework to conduct a Bayesian sweep across seven hyperparameters. The search space included network depth (conv blocks), filter width, dropout rates, and an explicit \textit{augmentation factor} representing the severity of noise injected during training.

\section{Part C: Results and Analysis}
\label{sec:results}

\subsection{Experimental Findings}
Table \ref{tab:comparison} summarizes the performance of our Optuna-optimized robust model against a standard clean-targeted baseline.

\begin{table}[H]
\centering
\caption{Model Performance across Sim-to-Real Severity Levels (Means and StdDev over 4 seeds)}
\label{tab:comparison}
\begin{tabular}{@{}lll@{}}
\toprule
\textbf{Metric} & \textbf{Optimized Baseline} & \textbf{Optuna-Robust Model} \\ \midrule
Clean Accuracy & 49.60\% $\pm$ 0.61\% & 46.70\% $\pm$ 0.72\% \\
Low Noise ($0.1$) & 13.72\% $\pm$ 0.35\% & 48.29\% $\pm$ 0.58\% \\
Med Noise ($0.25$) & 3.20\% $\pm$ 0.20\% & 47.47\% $\pm$ 0.81\% \\
High Noise ($0.4$) & 1.71\% $\pm$ 0.08\% & 15.39\% $\pm$ 1.12\% \\ \bottomrule
\end{tabular}
\end{table}

\subsection{The Alignment Tax and Catastrophic Forgetfulness}
As shown in Figure \ref{fig:dropoff}, the baseline model suffers from catastrophic forgetting the moment any noise is introduced. Our robust model maintains stable performance through the ``Medium Noise'' tier, representing a 14x increase in field-readiness.

\begin{figure}[H]
    \centering
    % \includegraphics[width=0.8\textwidth]{outputs/final_accuracy_dropoff.png}
    \caption{Performance comparison: Robust vs Vanilla models showing the cross-over point where robustness overcomes peak clean accuracy.}
    \label{fig:dropoff}
\end{figure}

\section{Conclusion}
This study confirms that Bayesian Domain Randomization is an essential component of robotic vision systems. While a minor accuracy penalty is paid on idealized simulation data, the resulting model is dramatically more capable of surviving the unpredictable visual noise of the real world.

\bibliographystyle{plain}
\bibliography{references}

\end{document}
